% Page 023 — page imprimée 19
% C'était aussi ça la Révolution à Meyrueis !!


{\fontsize{10.15}{11.9}\selectfont

\begin{center}
{\Large\bfseries C’était aussi ça la Révolution à Meyrueis !}

\vspace{0.65em}

{\bfseries
La société populaire de Meyrueis\\
ou\\
L’apprentissage de la Démocratie\\
Ventôse An II à Prairial An III\\
Février 1794 -- Mai 1795
}
\end{center}

\vspace{0.45em}

Après la Révolution des centaines de « Sociétés Populaires » ont vu le jour dans beaucoup de villes et villages. Elles permettaient pour la première fois à des représentants du peuple de faire entendre leurs voix.



On pouvait penser que la révolution, avec l’abandon des privilèges, avait apporté, au moins dans nos villages de la France profonde, la paix, la liberté religieuse pour les protestants, et une période de concorde et d’espoir. Ce n’est pas évident.
Un registre de la Société Populaire de Meyrueis comportant 108 pages nous permet de mieux comprendre ce qui s’est passé du 5 Ventôse An II jusqu’au 11-12 Prairial An III. (1794-1795).

On s’aperçoit que les luttes politiques, les Girondins (partisans d’un fédéralisme orientés plutôt au centre) contre les Montagnards (l’extrême gauche avec les sans-culottes), par exemple n’étaient pas ignorées, loin s’en faut, dans les campagnes où l’on apprenait assez vite ce qui se passait à Paris : les altercations, les violences, les manifestation et les exécutions.
C’est ce qui apparaît clairement dans les comptes-rendus des séances de « cette société des amis de l’ordre, de la justice et de l’humanité » (sic).

J’ai relevé quelques aspects de leur déroulement. Les séances s’ouvraient souvent par des couplets patriotiques \textit{« par les cris chéris : « Vive la République une et indivisible ou la mort ! » ou encore. « Vive la République, vive la Convention Nationale ! ».}

\medskip

\textbf{Les Fidel Castro ou Mao Tsé Toung n’ont rien inventé !!!}

On votait des motions envoyées à la Convention nationale : ainsi le 15 ventôse an II elle écrit:

\medskip

--- La Société a fait don de charpie et de paires de bas. \textit{(pour les soldats en opérations).}

--- Les représentants du peuple ont exercé un droit d’humanité en ordonnant l’abolition de l’esclavage.

--- Elle souhaite le prompt envoi du Code Civil, qui \textit{« doit anéantir l’hydre de la chicane... ».}

Un peu plus tard le 4 germinal (4 mars 1794) lors de la visite du citoyen Bon représentant du peuple : \textit{« Citoyens, je me suis aperçu en entrant dans cette société que, parmi les citoyennes qui la composent et qui doivent en faire l’ornement, il en est quelques une qui portent à leur col des marques superstitieuses qu’on appelait croix, ce signe qui alluma la guerre civile et fit s’entr’égorger nos pères. Je demande que ces marques soient déposées sur le bureau et échangées en bonnet de la liberté (le bonnet phrygien) ».}

Et, comme les citoyennes ne marquent point d’empressement à les déposer, le président nomme le citoyen Dupont pour en faire le recensement...

\medskip

\textbf{Déjà en 1794 les signes ostentatoires étaient prohibés !!!}

}

